% Figure 2. Distinct Floquet mechanisms from committed 60-digit flanks.
\begin{figure*}[t]
\centering
\begin{tikzpicture}[x=1cm,y=1cm,line cap=round,line join=round,
  >=Stealth,every node/.style={font=\footnotesize}]
  % ============================================================
  % (a) physical Sp(4) trace plane
  % a in [-4.5,4.5], b in [-2.5,6.5]
  % x = 0.58(a+4.5), y = 0.38(b+2.5)
  % ============================================================
  \begin{scope}
    \node[font=\small,anchor=north west] at (-0.25,4.05) {(a)};
    \draw[line width=0.75pt] (0,0) rectangle (5.22,3.42);
    \foreach \x/\lab in {0.29/-4,2.61/0,4.93/4}
      \draw[line width=0.55pt] (\x,-0.08) -- (\x,0)
        node[below=1pt] {$\lab$};
    \foreach \y/\lab in {0.19/-2,1.71/2,3.23/6}
      \draw[line width=0.55pt] (-0.08,\y) -- (0,\y)
        node[left=1pt] {$\lab$};
    \node[below=8pt,font=\small] at (2.61,0) {$a$};
    \node[rotate=90,anchor=south,font=\small] at (-0.66,1.71) {$b$};

    % Delta=0: b=a^2/4+2
    \draw[line width=1.0pt,dash dot]
      plot[smooth] coordinates {
        (0.29,3.23) (0.87,2.565) (1.45,2.09) (2.03,1.805)
        (2.61,1.71) (3.19,1.805) (3.77,2.09) (4.35,2.565)
        (4.93,3.23)
      };
    % G-=0: b=-2a-2 from (-4,6) to (0,-2)
    \draw[line width=1.0pt,dashed] (0.29,3.23) -- (2.61,0.19);
    % G+=0: b=2a-2 from (0,-2) to (4,6)
    \draw[line width=1.15pt] (2.61,0.19) -- (4.93,3.23);
    \node[font=\scriptsize,anchor=west] at (4.12,2.42) {$G_+=0$};
    \node[font=\scriptsize,anchor=east] at (1.18,2.48) {$G_-=0$};
    \node[font=\scriptsize,anchor=west] at (3.73,2.28) {$\Delta=0$};

    % mixed vertex (a,b)=(0,-2)
    \fill[black] (2.61,0.19) circle (1.8pt);
    \node[font=\scriptsize,anchor=north west] at (2.69,0.16) {$(0,-2)$};

    % +1 event: a=1.7065, b=1.4145 -> x=3.5998, y=1.4875
    \fill[black] (3.600,1.488) circle (2.4pt);
    \node[font=\scriptsize,anchor=east] at (3.48,1.73) {$+1$};

    % HH event: a=1.7007, b=2.7242 -> x=3.5964, y=1.9852
    \fill[black] (3.596,1.985) -- ++(0.00,0.16)
      -- ++(0.12,-0.24) -- ++(-0.24,0) -- cycle;
    \node[font=\scriptsize,anchor=west] at (3.76,2.12) {HH};
  \end{scope}

  % ============================================================
  % (b) zoom near +1
  % Re lambda in [0.94,1.08], Im lambda in [-0.06,0.06]
  % ============================================================
  \begin{scope}[shift={(6.15,0.05)}]
    \node[font=\small,anchor=north west] at (-0.20,4.00) {(b)};
    \draw[line width=0.75pt] (0,0) rectangle (3.35,3.35);
    \foreach \x/\lab in {0.000/0.94,1.436/1.00,2.871/1.06}
      \draw[line width=0.55pt] (\x,-0.08) -- (\x,0)
        node[below=1pt,font=\scriptsize] {$\lab$};
    \foreach \y/\lab in {0.000/-0.06,1.675/0,3.350/0.06}
      \draw[line width=0.55pt] (-0.08,\y) -- (0,\y)
        node[left=1pt,font=\scriptsize] {$\lab$};
    \node[below=8pt,font=\scriptsize] at (1.67,0)
      {$\operatorname{Re}\lambda$};
    \node[rotate=90,anchor=south,font=\scriptsize] at (-0.72,1.67)
      {$\operatorname{Im}\lambda$};
    \draw[line width=0.9pt,black!65]
      plot[domain=-3.5:3.5,samples=60]
        ({(cos(\x)-0.94)/0.14*3.35},{(sin(\x)+0.06)/0.12*3.35});
    \draw[line width=0.8pt] (1.436,1.675) circle (2.1pt);
    \fill[white] (1.436,1.675) circle (1.2pt);
    \node[font=\scriptsize,anchor=south] at (1.436,1.92) {$+1$};

    % unstable real reciprocal pair
    \draw[line width=0.9pt,fill=white]
      (0.833,1.675) ++(-0.11,-0.11) rectangle ++(0.22,0.22);
    \draw[line width=0.9pt,fill=white]
      (2.053,1.675) ++(-0.11,-0.11) rectangle ++(0.22,0.22);
    \node[font=\scriptsize,anchor=north] at (0.83,1.47) {real pair};

    % stable unit-circle pair
    \fill[black] (1.428,2.386) circle (2.4pt);
    \fill[black] (1.428,0.964) circle (2.4pt);
    \node[font=\scriptsize,anchor=west] at (1.62,2.55) {elliptic pair};

    \draw[->,line width=0.65pt] (0.98,1.75) -- (1.31,1.72);
    \draw[->,line width=0.65pt] (1.58,1.88) -- (1.52,2.20);
  \end{scope}

  % ============================================================
  % (c) zoom of opposite-Krein collision, upper half-plane
  % Re in [0.38,0.46], Im in [0.86,0.94]
  % ============================================================
  \begin{scope}[shift={(10.55,0.05)}]
    \node[font=\small,anchor=north west] at (-0.20,4.00) {(c)};
    \draw[line width=0.75pt] (0,0) rectangle (3.35,3.35);
    \foreach \x/\lab in {0.000/0.38,1.675/0.42,3.350/0.46}
      \draw[line width=0.55pt] (\x,-0.08) -- (\x,0)
        node[below=1pt,font=\scriptsize] {$\lab$};
    \foreach \y/\lab in {0.000/0.86,1.675/0.90,3.350/0.94}
      \draw[line width=0.55pt] (-0.08,\y) -- (0,\y)
        node[left=1pt,font=\scriptsize] {$\lab$};
    \node[below=8pt,font=\scriptsize] at (1.67,0)
      {$\operatorname{Re}\lambda$};
    \node[rotate=90,anchor=south,font=\scriptsize] at (-0.72,1.67)
      {$\operatorname{Im}\lambda$};
    \draw[line width=0.9pt,black!65]
      plot[domain=61:68,samples=40]
        ({(cos(\x)-0.38)/0.08*3.35},{(sin(\x)-0.86)/0.08*3.35});

    % stable modes and Krein signs
    \fill[black] (1.204,2.204) circle (2.6pt);
    \node[font=\scriptsize,anchor=east] at (1.02,2.35) {$+$};
    \draw[line width=0.95pt,fill=white] (2.580,1.558) circle (2.4pt);
    \node[font=\scriptsize,anchor=west] at (2.78,1.70) {$-$};

    % upper-half representatives of unstable quartet
    \draw[line width=0.9pt,fill=white]
      (1.568,1.209) ++(-0.11,-0.11) rectangle ++(0.22,0.22);
    \fill[black] (2.214,2.584) -- ++(0.00,0.14)
      -- ++(0.12,-0.22) -- ++(-0.24,0) -- cycle;
    \node[font=\scriptsize,align=center] at (1.93,0.67)
      {off-unit-circle\\quartet representatives};
  \end{scope}
\end{tikzpicture}
\caption{\label{fig:floquet}Distinct Floquet mechanisms on the continuation
component. (a)~Stability geometry in the physical $\operatorname{Sp}(4)$ trace
plane. The boundaries $G_+=0$ (solid), $G_-=0$ (dashed), and $\Delta=0$
(dot-dashed) correspond to $+1$, $-1$, and mode-collision conditions. The
filled circle marks the lower $+1$ transition and the filled triangle the upper
Hamiltonian--Hopf transition; the mixed vertex is $(a,b)=(0,-2)$.
(b)~Near the lower boundary, a real reciprocal pair (open squares) reaches
$+1$ and becomes an elliptic unit-circle pair (filled circles).
(c)~On the stable flank of the upper boundary, two unit-circle modes have
opposite Krein signatures (filled $+$ and open $-$). Beyond their collision,
the upper-half representatives lie on opposite sides of the unit circle; their
conjugates complete a reciprocal complex quartet.}
\end{figure*}
